\subsection{IDEALS}

DEFINITION 4. Let A be a ring. A subset \(\mathfrak{a}\) of A is called a left (resp. right) ideal if it is a subgroup of the additive group of A and the relations \(a \in \mathbf{A}\) , \(x \in \mathfrak{a}\) imply \(ax \in \mathfrak{a}\) (resp. \(xa \in \mathfrak{a}\) ). \(\mathfrak{a}\) is called a two-sided ideal of A if it is both a left ideal and a right ideal of A.

The definition of a left ideal may be expressed by the relations

\[
0 \in a, \quad a + a \subset a, \quad A. a \subset a
\]

the relation \(-a \subset a\) following from the formula \((-1).x = -x\) and \(A.a \subset a\) . For all \(x \in A\) , let \(\gamma_x\) be the mapping \(a \mapsto xa\) of \(A\) into \(A\) ; the action \(x \mapsto \gamma_x\) gives the additive group \(A^+\) of \(A\) the structure of a group with operators with \(A\) as set of operators. The left ideals of \(A\) are just the subgroups of \(A^+\) which are stable under this action.

The left ideals in the ring A are just the right ideals of the opposite ring \(A^{0}\) . In a commutative ring the three species of ideals are the same; they are simply called ideals.

Examples. (1) Let A be a ring. The set A is a two-sided ideal of A; so is the set consisting of 0, which is called the zero ideal and sometimes denoted by 0 or (0) instead of \{0\}.

\begin{enumerate}
\def\labelenumi{(\arabic{enumi})}
\setcounter{enumi}{1}
\item
  For every element \(a\) of \(\mathbf{A}\) , the set \(\mathbf{A}.a\) of left multiples of \(a\) is a left ideal; similarly the set \(a.\mathbf{A}\) is a right ideal. When \(a\) is in the centre of \(\mathbf{A}\) , \(\mathbf{A}.a = a.\mathbf{A}\) ; this ideal is called the principal ideal generated by \(a\) and is denoted by (a). (a) = A if and only if \(a\) is invertible.
\item
  Let M be a subset of A. The set of elements \(x \in A\) such that xy = 0 for all \(y \in M\) is a left ideal of A called the left annihilator of M. The right annihilator of M is defined similarly.
\item
  Every intersection of left (resp. right, two-sided) of A is a left (resp. right, two-sided) ideal. Given a subset X of A, there thus exists a smallest left (resp. right, two-sided) ideal containing X; it is called the left (resp. right, two-sided) ideal generated by X.
\end{enumerate}

Let \(\mathfrak{a}\) be a left ideal of \(\mathbf{A}\) . The conditions \(1 \notin \mathfrak{a}, \mathfrak{a} \neq \mathbf{A}\) are obviously equivalent.

DEFINITION 5. Let A be a ring. By an abuse of language, a left ideal a is said to be maximal if it is a maximal element of the set of left ideals distinct from A.

In other words, a is maximal if \(a \neq A\) and the only left ideals of A containing a are a and A.

THEOREM 1 (Krull). Let A be a ring and a a left ideal of A distinct from A. There exists a maximal ideal m of A containing a.

Consider A as operating on the additive group \(A^{+}\) of A by left multiplication. Then the left ideals of A are the stable subgroups of \(A^{+}\) . The theorem thus follows from §4, no. 3, Proposition 3 applied to the subset \(P = \{1\}\) of \(A^{+}\) .

PROPOSITION 3. Let A be a ring, \((x_{\lambda})_{\lambda \in L}\) a family of elements of A and a (resp. b) the set of sums \(\sum_{\lambda \in L} a_{\lambda} x_{\lambda}\) where \((a_{\lambda})_{\lambda \in L}\) is a family with finite support of elements of A (resp. \(\sum_{\lambda \in L} a_{\lambda} x_{\lambda} b_{\lambda}\) where \((a_{\lambda})_{\lambda \in L}, (b_{\lambda})_{\lambda \in L}\) are families with finite support of elements of A). Then a (resp. b) is the left (resp. two-sided) ideal of A generated by the elements \(x_{\lambda}\) .

The formulae

\[
0 = \sum_ {\lambda \in L} 0. x _ {\lambda}\tag{18}
\]

\[
\sum_ {\lambda \in L} a _ {\lambda} x _ {\lambda} + \sum_ {\lambda \in L} a _ {\lambda} ^ {\prime} x _ {\lambda} = \sum_ {\lambda \in L} (a _ {\lambda} + a _ {\lambda} ^ {\prime}) x _ {\lambda}\tag{19}
\]

\[
a. \sum_ {\lambda \in L} a _ {\lambda} x _ {\lambda} = \sum_ {\lambda \in L} (a a _ {\lambda}) x _ {\lambda}\tag{20}
\]

prove that \(\mathfrak{a}\) is a left ideal. Let \(\mathfrak{a}'\) be a left ideal such that \(x_{\lambda} \in \mathfrak{a}'\) for all \(\lambda \in \mathbf{L}\) and let \((a_{\lambda})_{\lambda \in L}\) be a family with finite support in A. Then \(a_{\lambda}x_{\lambda}\in \mathfrak{a}'\) for all \(\lambda \in L\) , whence \(\sum_{\lambda \in L}a_{\lambda}x_{\lambda}\in \mathfrak{a}'\) ; hence \(\mathfrak{a}\subset \mathfrak{a}'\) . Hence \(\mathfrak{a}\) is the left ideal of A generated by the \(x_{\lambda}\) . The argument for \(\mathfrak{b}\) is analogous.

PROPOSITION 4. Let A be a ring and \((\mathfrak{a}_{\lambda})_{\lambda \in L}\) a family of left ideals of A. The left ideal generated by \(\bigcup_{\lambda \in L} \mathfrak{a}_{\lambda}\) consists of the sums \(\sum_{\lambda \in L} y_{\lambda}\) where \((y_{\lambda})_{\lambda \in L}\) is a family with finite support such that \(y_{\lambda} \in \mathfrak{a}_{\lambda}\) for all \(\lambda \in L\) .

Let \(\mathfrak{a}\) be the set of sums \(\sum_{\lambda \in \mathbf{L}} y_{\lambda}\) with \(y_{\lambda} \in \mathfrak{a}_{\lambda}\) for all \(\lambda \in \mathbf{L}\) . The formulae \(\sum_{\lambda \in \mathbf{L}} x_{\lambda} + \sum_{\lambda \in \mathbf{L}} y_{\lambda} = \sum_{\lambda \in \mathbf{L}} (x_{\lambda} + y_{\lambda})\) and \(a\) . \(\sum_{\lambda \in \mathbf{L}} x_{\lambda} = \sum_{\lambda \in \mathbf{L}} ax_{\lambda}\) shows that \(\mathfrak{a}\) is a left ideal of \(\mathbf{A}\) . Let \(\lambda \in \mathbf{L}\) and \(x \in \mathfrak{a}_{\lambda}\) ; write \(y_{\lambda} = x\) and \(y_{\mu} = 0\) for \(\mu \neq \lambda\) ; then \(x = \sum_{\lambda \in \mathbf{L}} y_{\lambda}\) , whence \(x \in \mathfrak{a}\) and finally \(\mathfrak{a}_{\lambda} \subset \mathfrak{a}\) . If a left ideal \(\mathfrak{a}'\) contains \(\mathfrak{a}_{\lambda}\) for all \(\lambda \in \mathbf{L}\) , it obviously contains \(\mathfrak{a}\) and hence \(\mathfrak{a}\) is generated by \(\bigcup_{\lambda \in \mathbf{L}} \mathfrak{a}_{\lambda}\) .

The ideal \(\mathfrak{a}\) generated by \(\bigcup_{\lambda \in L} \mathfrak{a}_{\lambda}\) is called the sum of the left ideals \(\mathfrak{a}_{\lambda}\) and is denoted by \(\sum_{\lambda \in L} \mathfrak{a}_{\lambda}\) (cf.~II, § 1, no. 7). In particular, the sum \(\mathfrak{a}_{1} + \mathfrak{a}_{2}\) of the two left ideals consists of the sums \(a_{1} + a_{2}\) where \(a_{1} \in \mathfrak{a}_{1}\) and \(a_{2} \in \mathfrak{a}_{2}\) .

